\section{Solusi Distribusional untuk Persamaan Diferensial Biasa}

\begin{notn} Misalkan $\Omega$ suatu subhimpunan terbuka tak kosong dari $\R$, $n \in \N$, dan $a_0$, $a_1$, \dots $a_n \in \fml C^\infty(\Omega)$.
Kita meninjau operator diferensial (biasa)
 \index{diferensial!operator}%
 \index{operator!diferensial}%
berikut:
   \[ L = \sum_{k=0}^n a_k(x)\frac{d^k}{dx^k} \]
beserta persamaan diferensial (biasa)
 \index{diferensial!persamaan}%
 \index{persamaan!diferensial}%
yang terkait dengannya,
   \begin{equation}\label{eqn_ODE1}
        Lu = v
   \end{equation}
dengan $u$ dan $v$ distribusi-distribusi pada~$\Omega$. Sehubungan dengan Persamaan~\eqref{eqn_ODE1}, kita juga dapat meninjau dua variasi:
persamaan
   \begin{equation}\label{eqn_ODE2}
        \langle Lu, \phi\rangle = \langle v, \phi\rangle
   \end{equation}
dengan $\phi$ suatu fungsi uji, dan persamaan
   \begin{equation}\label{eqn_ODE3}
        \langle u, L^*\phi\rangle = \langle v, \phi\rangle
   \end{equation}
dengan $\phi$ suatu fungsi uji dan $L^* := \displaystyle\sum_{k=0}^n (-1)^k \frac{d^k(a_k\phi)}{dx^k}$ merupakan
 \index{adjoin!formal}%
 \index{adjoin formal}%
\df{adjoin formal} dari~$L$.
\end{notn}

\begin{defn} Andaikan bahwa dalam persamaan-persamaan di atas distribusi $v$ diberikan. Jika terdapat distribusi regular $u$ yang
memenuhi~\eqref{eqn_ODE1}, kita mengatakan bahwa $u$ merupakan solusi
 \index{solusi klasik}%
 \index{solusi!klasik}%
\df{klasik} bagi persamaan diferensial itu. Jika $u$ merupakan distribusi regular yang bersesuaian dengan suatu fungsi yang
\emph{tidak} memenuhi~\eqref{eqn_ODE1} dalam pengertian klasik, tetapi \emph{memenuhi}~\eqref{eqn_ODE2} untuk setiap fungsi uji~$\phi$,
maka kita mengatakan bahwa $u$ merupakan solusi
 \index{lemah!solusi}%
 \index{solusi!lemah}%
\df{lemah} bagi persamaan diferensial itu. Selanjutnya, jika $u$ merupakan distribusi singular yang memenuhi~\eqref{eqn_ODE3} untuk setiap fungsi
uji~$\phi$, kita mengatakan bahwa $u$ merupakan solusi
 \index{solusi distribusional}%
 \index{solusi!distribusional}%
\df{distribusional} bagi persamaan diferensial itu. Keluarga solusi
 \index{solusi diperumum}%
 \index{solusi!diperumum}%
\df{diperumum} mencakup semua solusi klasik, lemah, dan distribusional.
\end{defn}

\begin{exam} Solusi diperumum bagi persamaan $\dfrac{du}{dx} = 0$ pada~$\R$ hanyalah distribusi konstan (yakni,
distribusi regular yang berasal dari fungsi konstan). Jadi, setiap solusi diperumum merupakan solusi klasik.
\end{exam}

\begin{proof}[\emph{Petunjuk untuk bukti}] Andaikan $u$ suatu solusi diperumum bagi persamaan itu. Pilih suatu fungsi uji $\sbsb\phi 0$ sedemikian
sehingga $\int_\R \sbsb\phi 0 = 1$. Misalkan $\phi$ suatu fungsi uji sebarang pada~$\R$. Definisikan
   \[ \psi(x) = \phi(x) - \sbsb\phi 0(x)\int_\R \phi \]
untuk semua $x \in \R$. Amati bahwa $\langle u,\psi \rangle = 0$, lalu hitung $\langle u,\phi \rangle$.  \ns
\end{proof}

\begin{exam} Distribusi Heaviside adalah solusi lemah persamaan diferensial $x\dfrac{du}{dx}~=~0$ pada~$\R$. Tentu saja,
distribusi-distribusi konstan merupakan solusi klasik.
\end{exam}

\begin{exam} Distribusi Dirac $\delta$ merupakan solusi distribusional bagi persamaan diferensial $x^2\dfrac{du}{dx} = 0$ pada~$\R$.
(Distribusi Heaviside merupakan solusi lemah; dan distribusi-distribusi konstan merupakan solusi klasik.)
\end{exam}

\begin{exam} Untuk $p = 0,1,2,\dots$, turunan ke-$p$, $\dfrac{d^p\delta}{dx^p}$, dari distribusi delta Dirac merupakan solusi
bagi persamaan diferensial $x^{p+2}\dfrac{du}{dx} = 0$ pada~$\R$.
\end{exam}

\begin{exer} Carilah suatu solusi distribusional bagi persamaan diferensial $x \dfrac{du}{dx} + u = 0$ pada~$\R$.
\end{exer}










